\section{\bt{Introduction}{引言}}
\label{sec:introduction}
\input{figures/metric-record}
\bi{Large language models (LLMs) are enabling data agents for database analysis and predictive modeling~\cite{lei2025spider2,guo2024dsagent,liu2026agents}. These tasks rely on database-specific measures, grain, filters, and join paths~\cite{floratou2024nl2sql,robinson2024relbench}. For example, March, September, and third-quarter revenue share one definition (Figure~\ref{fig:metric-record}). Independent sessions repeat schema inspection, value probes, join tests, and the associated LLM calls. Across 100 memoryless metric sessions, 93\% of lookups target facts found by earlier sessions, and 71\% of session time precedes the final answer query.}
{大语言模型（LLM）正推动面向数据库分析与预测建模的数据智能体发展~\cite{lei2025spider2,guo2024dsagent,liu2026agents}。这些任务依赖数据库特有的度量、粒度、过滤条件与连接路径~\cite{floratou2024nl2sql,robinson2024relbench}。例如，3 月、9 月与第三季度的营业额共用同一定义（图~\ref{fig:metric-record}）。独立会话重复执行模式检查、取值探查、连接测试及相关的大模型调用。在我们追踪的 100 个无记忆指标会话中，93\% 的查找以此前会话已获得的事实为目标，71\% 的会话时间消耗在最终答案查询开始之前。}

\bi{Prior work supports reuse through agent memory and explicit database semantics. Agent memories retain insights, workflows, and playbooks~\cite{zhao2024expel,wang2025awm,zhang2026ace}, or share trajectories and composite tools~\cite{zhang2025gmemory,biswal2026agentsm}. Database-oriented systems recover calculation logic from scripts~\cite{weng2025datalab}, centralize metrics for SQL generation~\cite{dbt2026metricflow,databricks2026metricviews}, or construct predictive features through join-path and operator search~\cite{liang2025featpilot,fan2026deepprep}. Continued reuse also requires validating the computation's data assumptions.}
{既有工作通过智能体记忆与显式数据库语义支持复用。智能体记忆保留洞见、工作流与经验手册~\cite{zhao2024expel,wang2025awm,zhang2026ace}，或共享轨迹与复合工具~\cite{zhang2025gmemory,biswal2026agentsm}。面向数据库的系统从脚本恢复计算逻辑~\cite{weng2025datalab}，集中管理指标以生成 SQL~\cite{dbt2026metricflow,databricks2026metricviews}，或通过连接路径与算子搜索构造预测特征~\cite{liang2025featpilot,fan2026deepprep}。持续复用还要求验证计算所依赖的数据前提。}

\bi{Data changes can invalidate a definition without changing its SQL or schema. A revenue definition learned from one row per sale silently double-counts when a restatement retains historical and current rows. Static DSL syntax and structure checks~\cite{weng2025datalab} and unvalidated join declarations, such as Databricks' \code{at\_most\_one\_match}~\cite{databricks2026metricviews}, do not test grain or join-cardinality invariants. Version-based invalidation~\cite{wu2026invalidation} and derived-artifact repair~\cite{zhao2026memorepair} react to change; changed dependencies alone do not identify violated database assumptions. Sustained reuse thus requires executable definitions with checkable data conditions whose validation and upkeep are shared across agents.}
{数据变化能够在 SQL 与表模式均未改变的情况下使定义失效。从每笔销售占一行的数据中学到的营业额定义，在数据更正同时保留历史行与当前行后，会无报错地重复计数。领域专用语言（DSL）的语法与结构静态检查~\cite{weng2025datalab}，以及未经验证的连接声明（如 Databricks 的 \code{at\_most\_one\_match}~\cite{databricks2026metricviews}），都不检验粒度与连接基数不变式。基于版本的失效机制~\cite{wu2026invalidation}与派生产物修复~\cite{zhao2026memorepair}能够响应变化，但仅知道依赖发生变化，并不能确定哪些数据库前提已经失效。因此，持续复用需要为可执行定义附上可检验的数据条件，并在智能体之间共享其验证与维护。}

\bi{Our goal is to make discovered knowledge executable and safely reusable as databases evolve. This entails three challenges. \emph{First, extracting reusable computation.} Successful traces mix business logic, such as which sales rows qualify, with task parameters, such as a month or prediction cutoff. Generalization must preserve the former while parameterizing the latter. \emph{Second, establishing validity at use time.} Conditions must capture grain and join assumptions, with verification evidence valid on the consuming query's exact snapshot despite concurrent writes. \emph{Third, amortizing maintenance.} Revalidating definitions independently for each agent repeats costly scans of shared dependencies, while discarding them after every write forces agents back into costly re-exploration. Shared maintenance must eliminate redundant checks while preserving snapshot consistency.}
{我们的目标是使探索所得的知识可执行，并能在数据库演进过程中安全地复用。这面临三大挑战。\emph{第一，提取可复用的计算。}成功轨迹混合了业务逻辑（如哪些销售记录符合统计条件）与任务参数（如月份或预测截点）；泛化过程必须保留前者，并将后者参数化。\emph{第二，建立使用时的有效性。}条件必须刻画粒度与连接假设；即使存在并发写入，验证证据也必须对消费查询实际读取的快照成立。\emph{第三，分摊维护开销。}为每个智能体独立重验定义，会在共同依赖的数据上重复执行昂贵的扫描；而每次写入后丢弃定义，又会迫使智能体退回高代价的重复探索。共享维护必须在保证快照一致性的前提下消除冗余检查。}

\bi{We present \system\ (\textbf{M}etric-\textbf{A}ware \textbf{V}alidation and \textbf{R}euse for \textbf{A}gents), a shared memory system between data agents and the database. We first represent executable metrics and features by reusable measures, grain, joins, and filters, with entities and temporal windows for features. Successful sessions supply candidate definitions; deterministic publication checks validate their structure and compare compiled results with verified metric answers or submitted feature values. Compilation instantiates accepted definitions for new periods and prediction instances.}
{本文提出 \system{}（\textbf{M}etric-\textbf{A}ware \textbf{V}alidation and \textbf{R}euse for \textbf{A}gents），一个位于数据智能体与数据库之间的共享记忆系统。首先，我们以可复用的度量、粒度、连接和过滤表示可执行的指标与特征，并为特征记录实体与时间窗口。成功会话提供候选定义；确定性的发布检查验证其结构，并将编译结果与经过核实的指标答案或提交的特征值比较。编译过程为新的期间与预测样本实例化已通过检查的定义。}

\bi{We then develop shared validation and maintenance. Results indexed by condition identity and dependency versions let definitions and agents share validation work. Updates trigger revalidation of affected conditions; failures trigger checked metric repair or revocation. For metric queries, snapshot-bound execution verifies that all required conditions hold on the consuming snapshot; feature compilation enforces per-instance temporal bounds. Validated equivalent metric rewrites are also shared. We establish sufficient conditions for metric aggregation, sound check reuse, and rewrite equivalence.}
{随后，我们设计共享验证与维护机制。以条件标识和依赖版本为索引的检查结果，使不同定义与智能体能够共享验证工作。数据更新触发对受影响条件的重验；检查失败则触发经过验证的指标修复或撤销。对于指标查询，快照绑定执行机制验证所有必要条件在消费快照上成立；特征编译则落实逐样本的时间边界。经过验证的等价指标改写同样向使用者共享。本文给出指标聚合的充分条件、检查结果复用的可靠性依据及改写等价性证明。}

\bi{We evaluate \system\ on RelBench~\cite{robinson2024relbench}, TPC-DS, and TPC-H. On TPC-DS, \PpAgents\ agents use \PpMavraTokSaveNone\% fewer tokens and run \PpMavraSpeedupNone$\times$ faster than without memory, using \PpMavraTokSavePrivate\% fewer tokens than private memory with the same mechanisms. Post-restatement accuracy is \PpMavraChangeAcc\%, versus \PpNoneChangeAcc\% without memory. In the independent official-API RelBench study, reuse-stage tokens (cutoffs 2--3) fall by \OfStackSaveReuse\% on rel-stack and \OfTrialSaveReuse\% on rel-trial relative to no memory. Across all cutoffs, five of six tasks meet the predeclared quality margins; study-adverse has \OfAdverseRegression\% higher MAE. The full-sequence token savings are \OfStackSaveAll\% and \OfTrialSaveAll\%, including memory operations and failed attempts.}
{我们在 RelBench~\cite{robinson2024relbench}、TPC-DS 和 TPC-H 上评估 \system。在 TPC-DS 上，\PpAgents\ 个智能体相对无记忆减少 \PpMavraTokSaveNone\% 的 token，快 \PpMavraSpeedupNone\ 倍；相对采用相同机制的私有记忆减少 \PpMavraTokSavePrivate\% 的 token。数据更正后，正确率为 \PpMavraChangeAcc\%，无记忆组为 \PpNoneChangeAcc\%。在独立的官方API RelBench实验中，复用阶段（截点2--3）的token相对无记忆在rel-stack与rel-trial上分别减少\OfStackSaveReuse\%与\OfTrialSaveReuse\%。全部截点上六项任务中五项满足预设质量门槛，study-adverse的MAE高\OfAdverseRegression\%。计入记忆操作和失败尝试后，全序列token分别减少\OfStackSaveAll\%与\OfTrialSaveAll\%。}

\bi{Our contributions are:}{本文的主要贡献如下：}
\begin{itemize}
\item \bi{An executable knowledge representation coupling metric and feature definitions with data conditions and learning evidence for reuse across questions and agents.}{一种可执行知识表示，将指标与特征定义同数据条件和学习证据结合，支持跨问题与智能体复用。}
\item \bi{A shared validation and maintenance protocol binding metric validity to consuming snapshots and amortizing checks, repairs, and optimizations, with formal results for aggregation, check reuse, and rewrite equivalence.}{一套共享验证与维护协议，将指标有效性绑定到消费快照，分摊检查、修复与优化开销，并给出聚合、检查复用与改写等价性的形式化结果。}
\item \bi{A workload analysis, cost model, and evaluation on RelBench, TPC-DS, and TPC-H quantifying exploration redundancy, sharing benefits, maintenance costs, and correctness under updates.}{通过负载分析、代价模型及 RelBench、TPC-DS 和 TPC-H 上的实验，量化探索冗余、共享收益、维护代价与更新下的正确性。}
\end{itemize}
